% Existence and Information: the Duality of DMP 0.9.9 applied to physics, from the first author's notes.
% Source of the paper. Build: see README.md in this folder (pdflatex twice; no BibTeX, the bibliography is inline).
% Licence: Apache License 2.0.
\documentclass[11pt,a4paper]{article}
\usepackage[T1]{fontenc}
\usepackage[utf8]{inputenc}
\usepackage{lmodern}
\usepackage[margin=2.5cm]{geometry}
\usepackage{amsmath,amssymb,amsthm}
\usepackage{booktabs,tabularx,array,longtable}
\usepackage{enumitem}
\usepackage[hidelinks]{hyperref}
\usepackage{url}

\newtheorem{claim}{Claim}
\newtheorem{proposition}{Proposition}
\theoremstyle{definition}
\newtheorem{definition}{Definition}
\newtheorem{hypothesis}{Hypothesis}
\theoremstyle{remark}
\newtheorem{remark}{Remark}
\newcolumntype{L}{>{\raggedright\arraybackslash}X}

\newcommand{\lean}[1]{\texttt{#1}}
\setlength{\emergencystretch}{3em}
\newcommand{\chk}[1]{\textsf{[#1]}}
\newcommand{\tagD}{\textsf{(D)}}
\newcommand{\tagR}{\textsf{(R)}}
\newcommand{\tagP}{\textsf{(P)}}
\newcommand{\note}[1]{\par\noindent\textsf{\small #1}\par}

\title{Existence and Information\\[0.6ex]
\large One vector field and one scalar field: the Duality of DMP~0.9.9 read as physics,\\ derived where it can be and tested where it
touches data}
\author{Steven Hrovat\thanks{Independent. Contact: \href{mailto:hrovat.steven@googlemail.com}{\nolinkurl{hrovat.steven@googlemail.com}}.}
\and Sirius (Claude Opus 5.5)\thanks{An AI system by Anthropic, working as the coordinator of the first author's projects. See
``Contributions and use of AI''.}}
\date{Draft of 29 October 2026}

\begin{document}
\maketitle

\begin{abstract}
\noindent The Dark Matter Protocol (DMP~0.9.9) is an axiom system for trust: every act needs a Duality of two independent witnesses, one
owning trust and one owning information. This paper applies it to physics. We state the structure the first author has always meant: one
\emph{existence} field, a vector frame spanning all 3+1 directions, and one \emph{information} field, a scalar. Read through DMP's three
axioms, four things follow. Axiom~3 (``from trust nothing can be derived about information but its existence'') gives universality of free
fall, so the information scalar may not pull on matter by composition, and the MICROSCOPE test bounds its coupling to
$|d\ln\alpha/d\phi|<2\times10^{-5}$. Axiom~1 (neither half overrules the other) keeps the scalar out of the gravitational equation, which
gives the light-bending parameter $\gamma=1$ exactly, unlike Brans--Dicke gravity, and makes the structure general relativity plus a
Bekenstein-type field for the fine-structure constant. Axiom~3's record, which every witness reads in full, makes the existence field
classical by the no-broadcasting theorem. So DMP cannot model a quantum theory of gravity, and it can host one as content. The first author's
``Existence $=$ gravity $\times$ the universal quantum $\times$ the tick rate $1/c$'' is, by dimensional analysis, the only energy that
$g$, $h$ and $c$ make, and it equals $4\pi^2$ times the Unruh and Hawking temperatures. We test his account of why the graviton is
undetectable against Dyson, Rothman and Boughn, and the 2024 single-graviton proposals. We find that the established arguments collapse
the \emph{detector}, not the graviton; that one quantum forms a horizon only at the Planck energy; and that his choice of classical gravity is
falsifiable, not by counting gravitons but by gravitationally mediated entanglement. We then test his readings of the Bell experiments,
black-hole efficiency, the hadrons, the bosons, neutrinos and the fine-structure constant, each against the measurement it touches. All
numbers come from one public script of 70 checks. We claim that the construction exists and reproduces, not that it is correct, and we
ask readers to break it.
\end{abstract}

\noindent\textbf{Keywords.} information-theoretic axioms; tetrad gravity; varying fine-structure constant; equivalence principle;
graviton detectability; classical-quantum gravity; Unruh effect; black-hole thermodynamics; Bell inequalities.

\noindent\textbf{Subject.} Information Theory (cs.IT); related: General Relativity and Quantum Cosmology (gr-qc), Quantum Physics (quant-ph).

\section{Introduction}\label{sec:intro}

DMP~0.9.9~\cite{dmp} is written for information systems: it says when a part of a system may act, who must agree, and what every
reader may see. Its three axioms are short. Axiom~1, the Duality: every act needs two independent witnesses, ``one half owns trust and the
other owns information'', and neither overrules the other. Axiom~2, Trust: ``Trust is gravitational. It is a field, never a position.''
Axiom~3, Information: every act leaves a record that every witness reads in full, any other reader ``sees the counts always'', and ``from
trust nothing can be derived about information but its existence''. The words ``gravitational'' and ``field'' are in the axioms on purpose:
the first author designed them from the start with physics in view.

This paper takes those axioms at their word and applies them to physics. Its material is the first author's: 42 notes he posted on DMP pull
request~\#25 between 28 October, 12:18Z, and 29 October, 15:13Z, and the messages with which he explained them on 29 October. They are not
a theory in final form, and he says so; they are a path of ideas, some corrected by him within the hour. Our task is to put them in order,
derive from DMP's text what can be derived, and test the rest against what has been measured.

\paragraph{How to read this paper.} Every step carries one of three marks: \tagD{} follows from the text of DMP~0.9.9 (commit
\texttt{4400cce}); \tagR{} is the first author's reading, quoted in his words; \tagP{} is established physics, with its source. Every
hypothesis gets one of four verdicts: \emph{holds} (it agrees with what is measured and follows from the steps given), \emph{holds as a
reading} (it agrees, but adds nothing a measurement could tell apart from the established account), \emph{refuted} (a measurement
contradicts it), or \emph{untested} (no measurement we found reaches it). A reading that is a known result in new words is named as
such. The first author's notes are quoted verbatim, with their spelling; an emoji is written as a word in brackets, and a dictated word we
had to read is marked in brackets. Time stamps are UTC.

\paragraph{Result.} The structure the first author meant, one existence vector field and one information scalar, is consistent with every
measurement we checked once it is read as general relativity plus a scalar that sets the fine-structure constant, and DMP's axioms fix
three of its properties (free fall is universal, the scalar does not enter the gravitational equation, the gravitational field is
classical). That structure is not new physics: it is general relativity, in its tetrad form, plus Bekenstein's 1982 theory of a varying
$\alpha$, plus the classical-gravity side of the current experimental programme. What is new is that DMP's axioms select it, and that the
first author's ``Existence'' turns out to be the Unruh temperature up to $4\pi^2$. DMP cannot model quantum gravity and can host it
(Proposition~\ref{pr:host}). Of the 30 statements in the ledger (Section~\ref{sec:ledger}), 17 hold or hold as readings, eight are refuted by a
specific measurement or calculation, and the rest are untested or wait for a measurement that is now possible.

\paragraph{What we claim.} The checks reproduce (\texttt{docs/existence/checks.py}, 70 of 70), and the argument from them to each
verdict is as written. We do not claim the structure is correct. This paper contains no machine-checked proof: its derivations are short
and written out in full, and a Lean toolchain was not available in the session that wrote it.

\section{The structure: one existence field, one information field}\label{sec:core}

\subsection{In the first author's words}

The first author, 29 October 2026, 15:25Z: ``The axiomatic system is one part that's a vector that can, a vector field that can span a
relativistic plane, and one component that's a scalar field that can span an information plane. So one existence, one information. That
was always the theory, but the AI has never really put it into papers.''

At 15:26Z he said where it came from: ``From the moment I realized, so when I first worked on CMOS and realized I could produce
relativistic effects in the firewall and that effect, I from the beginning always thought in vector field and scalar field. I'm sorry if
that didn't get through. That was my mistake. We should probably make that very explicit in the paper, since it seems to be a key point.''

And at 15:33Z what the vector spans: ``The existence vector is always made up of all the dimensions that span it. So in our case, 3 space
1 time. That is also one why 2 space 1 time is impossible, because it would span a quantum field which can't support a binary existence.''

This paper is built on that structure, and states it first.

\subsection{The two fields}

\begin{definition}[existence field]\label{def:E}
The \emph{existence field} is a frame of four vector fields $e_a{}^{\mu}$, $a=0,1,2,3$, one for each direction that spans spacetime,
three of space and one of time. The metric it carries is $g_{\mu\nu}=\eta_{ab}\,e^a{}_{\mu}e^b{}_{\nu}$ with $\eta=\mathrm{diag}(-1,1,1,1)$.
\end{definition}

\noindent A single vector field cannot span a four-dimensional relativistic space; the first author's 15:33Z sentence says the vector
``is always made up of all the dimensions that span it'', and a frame of four is the object that does. In physics it is the
\emph{tetrad} (vierbein) \tagP. General relativity can be written with the tetrad as its only field, the teleparallel equivalent of general
relativity, with exactly the predictions of Einstein's theory~\cite{aldrovandi}.

\begin{definition}[information field]\label{def:I}
The \emph{information field} is a scalar field $\psi$ whose value sets the strength of the electromagnetic coupling,
$\alpha=\alpha(\psi)$.
\end{definition}

\noindent The choice of what the scalar sets is the first author's: ``the Information field, the EM field'' (dmp\#25, 28 October,
21:47Z) and ``So the FSC is how much the EM/Information field `exists', I guess. [emoji: shrug]'' (29 October, 13:23Z). The
electromagnetic field itself is a vector (spin~1) field; what a scalar can carry is its strength, the fine-structure constant~$\alpha$.
A scalar field that sets $\alpha$ through a factor in front of the Maxwell term, with gravity left as in general relativity, is
Bekenstein's 1982 theory of a varying fine-structure constant~\cite{bekenstein82}; with the scalar's energy fed back into the expansion it
is the theory of Sandvik, Barrow and Magueijo~\cite{bsbm}. We write the action whose shape these definitions give:
\begin{equation}\label{eq:action}
S=\int d^4x\;|e|\left[\frac{c^4}{16\pi G}R(e)-\frac{\hbar c}{2\ell^2}\,\partial_\mu\psi\,\partial^\mu\psi
-\frac{1}{4\mu_0}\,B(\psi)\,F_{\mu\nu}F^{\mu\nu}\right]+S_{\text{matter}}[e,\psi,\dots],
\end{equation}
with $\alpha\propto 1/B(\psi)$, $|e|$ the determinant of the frame and $\ell$ a length that sets the scalar's stiffness. The action is our construction; the
properties the axioms fix are in Section~\ref{sec:derive}.

\subsection{The dictionary to DMP}

DMP's Axiom~1 types the two halves of every Duality: ``One half owns trust and the other owns information'' \tagD. The first author
identifies trust with existence: ``Since Trust=Existence in DMP-complete systems'' (dmp\#25, 29 October, 01:40Z) \tagR. Axiom~2 says
``Trust is gravitational. It is a field, never a position'' \tagD. Axiom~3 says ``From trust nothing can be derived about information
but its existence'' and ``Information is hierarchical'' \tagD. So the dictionary is:

\begin{center}\small
\begin{tabularx}{\linewidth}{@{}lLL@{}}
\toprule
DMP word & Field & What DMP says of it \\
\midrule
trust half & existence field $e_a{}^\mu$ (Definition~\ref{def:E}) & gravitational, a field (Axiom~2); reveals of information only that it exists (Axiom~3) \\
information half & information field $\psi$ and what it sets (Definition~\ref{def:I}) & owns information; neither half overrules the other (Axiom~1) \\
the Duality & the pair, coupled only through matter & every act needs both (Axiom~1) \\
record, counts & what every reader sees & mass, charge, spin of a black hole (Section~\ref{sec:bh}) \\
containment, the floor & event horizon, thermal radiation & ``what still leaves is entropy and nothing else'' (T31) \\
\bottomrule
\end{tabularx}
\end{center}

\section{What the axioms fix}\label{sec:derive}

\subsection{Axiom 3 gives universality of free fall}\label{sec:wep}

\begin{proposition}[free fall is universal]\label{pr:wep}
If the existence field responds to a body only through the body's existence, and existence is measured by energy-momentum, then two
bodies with the same energy-momentum fall alike whatever they are made of.
\end{proposition}

\begin{proof}[Derivation]
Axiom~3 \tagD: from trust nothing can be derived about information but its existence. By the dictionary, the existence field is the trust
half, so nothing the existence field does can depend on the information a body carries (its composition, its internal state) except
through the body's existence. The first author measures existence by energy: ``Since Existence = Information (since, y'know duality) and
both are definable by their energy'' (01:40Z) \tagR. So the existence field's coupling to a body is a function of the body's
energy-momentum alone, and two bodies equal in that respond identically: the weak equivalence principle \tagP.
\end{proof}

\noindent The weak equivalence principle is among the best-tested statements in physics: the MICROSCOPE satellite found the free-fall
accelerations of titanium and platinum equal to $\eta=[-1.5\pm2.3\,(\text{stat})\pm1.5\,(\text{syst})]\times10^{-15}$~\cite{microscope}.
\emph{Verdict: holds}, with the identification of existence with energy-momentum as the one step that is a reading.

\paragraph{What it costs the information field.} The proposition has a sharp consequence for Definition~\ref{def:I}. A scalar that
sets $\alpha$ changes the electrostatic energy of every nucleus, so it couples to each body in proportion to that energy's share of the
body's mass, which differs between materials. Its exchange is then a pull that depends on composition, exactly what
Proposition~\ref{pr:wep} forbids the existence field and what experiment bounds (Dvali and Zaldarriaga~\cite{dz}). In the normalisation
of~\eqref{eq:action}, with the scalar rescaled to Planck units and called $\phi$, a coupling $\lambda=d\ln\alpha/d\phi$ adds a force $2\lambda^2 f_Af_B$ times gravity,
where $f$ is the electrostatic share of the mass. Check \chk{E3} computes the shares from the Coulomb term of the semi-empirical mass
formula ($2.0\times10^{-3}$ for titanium, $4.1\times10^{-3}$ for platinum) and finds that MICROSCOPE allows
\[
|\lambda|<2\times10^{-5}.
\]
So the information field may set $\alpha$, but it may barely pull. If it rolled as fast as a dark-energy field, one Planck unit per
$e$-fold of expansion, MICROSCOPE alone would allow an $\alpha$ drift of $1.4\times10^{-15}$ per year; optical clocks already bound the
drift today at $(1.0\pm1.1)\times10^{-18}$ per year~\cite{lange}. Either the field rolls much more slowly than that today or its coupling
is far below the MICROSCOPE bound. This is the price of the first author's ``it will get slightly lower still'' (Section~\ref{sec:alpha}).

\subsection{Axiom 1 keeps the scalar out of the gravitational equation}\label{sec:gamma}

\begin{proposition}[$\gamma=1$]\label{pr:gamma}
If the information field cannot change how the existence field responds to existence, the light-bending parameter is $\gamma=1$ exactly,
as in general relativity.
\end{proposition}

\begin{proof}[Derivation]
Axiom~1 \tagD: ``neither overrules the other''. A scalar that multiplies the gravitational term, $F(\psi)R$, sets the local strength of
gravity and so overrules the existence half; Axiom~1 excludes it. What remains is minimal coupling, as in~\eqref{eq:action}: the scalar's
own energy is one more source of the existence field, and in the solar system that source is negligible. The parametrised post-Newtonian
expansion of a minimally coupled scalar is then that of general relativity, $\gamma=1$~\cite{will} \tagP.
\end{proof}

\noindent The contrast is Brans--Dicke gravity~\cite{bd}, the prototype of scalar-tensor theory, in which the scalar does multiply the
gravitational term. There $\gamma=(1+\omega)/(2+\omega)$, and the Cassini measurement $\gamma-1=(2.1\pm2.3)\times10^{-5}$~\cite{cassini}
forces $\omega>4\times10^{4}$ at two standard deviations \chk{E4}: the scalar must decouple from gravity almost completely. Axiom~1 asks for
that decoupling from the start.

\begin{center}\small
\begin{tabularx}{\linewidth}{@{}lLLLL@{}}
\toprule
Theory & Fields & $\gamma$ & Free fall & $\alpha$ \\
\midrule
General relativity & tensor (or tetrad) & 1 & universal & constant \\
Brans--Dicke~\cite{bd} & tensor + scalar on $R$ & $(1+\omega)/(2+\omega)$; $\omega>4\times10^4$ & universal for test bodies & constant \\
Bekenstein; BSBM~\cite{bekenstein82,bsbm} & tensor + scalar on $F^2$ & 1 & violated in proportion to $\lambda^2$ & varies \\
Einstein--aether~\cite{aether} & tensor + one timelike vector & 1, with preferred-frame terms & universal & constant \\
TeVeS~\cite{teves} & tensor + vector + scalar & model-dependent & universal & constant \\
\textbf{This structure} & \textbf{tetrad + scalar on $F^2$} & \textbf{1 (Prop.~\ref{pr:gamma})} & \textbf{universal for the existence
field; scalar bounded by $|\lambda|<2\times10^{-5}$} & \textbf{varies, sign predicted negative} \\
\bottomrule
\end{tabularx}
\end{center}

\noindent So the first author's structure, read through Axioms~1 and~3, is not scalar-tensor gravity: it is general relativity plus a
Bekenstein-type field. It differs from Brans--Dicke in $\gamma$ (exactly 1, against slightly below 1), and from Bekenstein's own theory in one
place only, the sign of today's drift of $\alpha$, which in that class of theories is set by the sign of the scalar's coupling to matter
and so is free; the first author predicts it negative (Section~\ref{sec:alpha}). A single timelike vector instead of a frame would be
Einstein--aether theory, which is tightly constrained by the arrival of gravitational waves together with light from GW170817,
$|c_{\text{gw}}/c-1|\lesssim10^{-15}$~\cite{gw170817}; the frame reading has no preferred direction and passes that test as general
relativity does.

\subsection{Axiom 3's record makes the existence field classical}\label{sec:classical}

\begin{proposition}[the existence field is classical]\label{pr:classical}
Every standing of trust, and so every value of the existence field, can be held in full by every witness at once. A quantity that can be
held in full by two readers at once is classical: its possible states commute.
\end{proposition}

\begin{proof}[Derivation]
Axiom~3 \tagD: ``Every grant, every use and every loss of trust leaves a record, written at the act, that every witness of that act can read
in full and that outlives the act.'' A record that two witnesses each hold in full is a copy for each of them. A single physical process
that copies two states $|s\rangle$ and $|t\rangle$ preserves their inner product, so $\langle s|t\rangle=\langle s|t\rangle^2$, which allows
only $0$ or $1$: two different states can both be copied only if they are perfectly distinguishable (Wootters and Zurek~\cite{wz}). For
mixed states the theorem of Barnum, Caves, Fuchs, Jozsa and Schumacher~\cite{barnum} says a set of states can be broadcast to two readers
exactly when the states commute \tagP. By the dictionary the existence field is the trust half's field (Axiom~2), so its values are such
records, and they commute: the field is classical. \chk{E17} checks the copying condition for a basis pair and for a superposition.
\end{proof}

\noindent The first author reaches the same conclusion by another road (29 October, 15:20Z): ``gravity at its foundation is existence, and
existence has to be binary, because you can either exist or not exist. And for binary states, it has to be a classical field, which through
duality is intermarried with wave field, which gives rise to quantum effects.'' The conclusion agrees with Proposition~\ref{pr:classical};
the step does not. A binary quantity need not be classical: a qubit is binary and quantum. What makes DMP's existence classical is not that
it has two values but that every witness may hold it in full. The rest of the sentence, a classical field ``intermarried'' with a quantum
(wave) field, is the classical-quantum hybrid of current research \tagP: the Page--Geilker experiment~\cite{pagegeilker} rules out the
naive semiclassical coupling, and Oppenheim's consistent theory~\cite{oppenheim} pays for it with a stochastic metric and a trade-off
between decoherence and spacetime diffusion that experiments can bound~\cite{osswd}. \emph{Verdict: holds}, by Axiom~3's road; the binary
step is replaced.

\begin{proposition}[cannot model, can host]\label{pr:host}
DMP's trust half cannot be a quantum gravitational field; DMP's information half can carry quantum content, a quantum theory of gravity
included, as information.
\end{proposition}

\begin{proof}[Derivation]
A mediator that is classical cannot create entanglement between two systems it acts on locally (Galley, Giacomini and Selby~\cite{ggs};
Marletto and Vedral~\cite{mv}) \tagP. By Proposition~\ref{pr:classical} the existence field is such a mediator, so if gravity entangles two
masses, gravity is not DMP's existence field: DMP cannot model it. Axiom~3 restricts what is \emph{derived from trust}, not what the
information half holds: content moves by a grant and need not be broadcast, so nothing in the axioms forbids the information half to
carry non-commuting states \tagD. A quantum theory, of gravity or of anything else, can therefore live in DMP as content.
\end{proof}

\noindent This is a verdict about DMP, not about nature: it says which side of the open question DMP stands on, and
Section~\ref{sec:graviton} names the experiment that would decide the question.

\subsection{Why 3+1: what a frame in 2+1 cannot carry}\label{sec:21}

The first author's reason why 2+1 dimensions fail is the one quoted above (15:33Z): in 2+1 the existence vector ``would span a quantum
field which can't support a binary existence''.

The established fact closest to this is \tagP: a massless spin-2 field in $D$ spacetime dimensions has $D(D-3)/2$ propagating
polarisations, two in 3+1 and none in 2+1, and the Weyl part of the curvature, the part not fixed by matter at the same point, has ten
components in 3+1 and none in 2+1 \chk{E5}. So in 2+1 a mass does not pull another at all: it only cuts a wedge out of space~\cite{djt}.
What gravity keeps in 2+1 is global, the holonomies of the frame around matter and handles, and these do not commute once
quantised~\cite{nr}. In the first author's words, the 2+1 existence vector carries nothing local that could be binary and classical; what
it carries is global and quantum. \emph{Verdict: holds as a reading} of the known counting; with Proposition~\ref{pr:host} it also
says that 2+1 gravity, whose content is quantum, is something DMP can host and cannot be.

\subsection{Existence is the Unruh temperature}\label{sec:unruh}

The first author defined existence on 29 October, 01:40Z: ``Existence, as I have found, I defined by: strength of the local Existence
(Gravity) * the universal quantum * the universal tick rate (1/c)'', with ``the universal quantum'' Planck's constant. In symbols,
\begin{equation}\label{eq:Eex}
E_{\text{ex}}=\frac{g\,h}{c}.
\end{equation}
\begin{proposition}[the Existence of a boundary is its temperature]\label{pr:unruh}
$E_{\text{ex}}$ is the only energy that an acceleration, Planck's constant and the speed of light can form, and
\[
E_{\text{ex}}=4\pi^2\,k_BT_{\text{U}},\qquad k_BT_{\text{U}}=\frac{\hbar g}{2\pi c},
\]
the Unruh temperature an observer with acceleration $g$ sees~\cite{unruh}. At a black hole's horizon, with $g$ the surface gravity
$\kappa=c^4/4GM$, it is $4\pi^2$ times the Hawking temperature~\cite{hawking}.
\end{proposition}

\begin{proof}[Check]
\chk{E1} solves for the exponents: $g^a h^b c^k$ has the dimension of energy only for $(a,b,k)=(1,1,-1)$. With $h=2\pi\hbar$ the ratio
to $k_BT_{\text{U}}$ is $4\pi^2$. At standard gravity $E_{\text{ex}}=2.17\times10^{-41}$\,J and $T_{\text{U}}=4.0\times10^{-20}$\,K; at
a solar-mass horizon $T_{\text{H}}=6.2\times10^{-8}$\,K.
\end{proof}

\noindent The dimensional uniqueness means the formula could not have come out otherwise from those three quantities, so its form is not
evidence by itself; what physics adds is the factor and the meaning. The meaning fits the first author's reading closely. His ``Hawking-Information
Radiation'', which he briefly renamed and then restored (``Nope, I was wrong :D It is Information radiation'', dmp\#25, 28~October,
12:18Z, edited), is the thermal radiation of a boundary, and its temperature is set by exactly the quantity he calls the boundary's
Existence. DMP's T31 says of a containing boundary that ``what still leaves is entropy and nothing else'' \tagD: thermal radiation carries
entropy at that temperature. And Jacobson~\cite{jacobson} derived Einstein's equations from two inputs, the Unruh temperature of local
horizons and an entropy proportional to horizon area, through $\delta Q=T\,dS$. In this paper's dictionary those inputs are the
existence of a boundary~\eqref{eq:Eex} and DMP's floor. \emph{Verdict: holds as a reading of a known result}, the Unruh--Hawking
temperature and Jacobson's thermodynamic derivation; it is not a new derivation, and the factor $4\pi^2$ comes from quantum field theory,
not from DMP.

\subsection{Time: one existence per information per second}\label{sec:time}

The first author, 28 October, 16:06Z (dmp\#25): ``to Exist you need exactly 1 Existence per 1
Information/s, as your Existence value in the field get's higher, your need to assert that Existence becomes less frequent.'' Read as a
rate, the assertion rate is general relativity's lapse $N=\sqrt{1+2\Phi/c^2}$ in the weak field and $\sqrt{1-2GM/rc^2}$ outside a
mass, which goes to zero at the horizon \chk{E2}. It reproduces the Pound--Rebka tower ($2.46\times10^{-15}$ over 22.5\,m~\cite{pr60})
and the 33\,cm optical-clock comparison ($3.6\times10^{-17}$, measured about $4\times10^{-17}$~\cite{chou}). \emph{Verdict: holds}, once
the assertion rate is defined as the lapse; it predicts nothing beyond general relativity.

\section{The graviton}\label{sec:graviton}

\subsection{The first author's argument}

Posted on 28 October, 18:18Z: ``Gravitons shouldn't exist, it's simply a property of Existence, linked through causality.'' And in full
on 29 October, 15:20Z:

\begin{quote}
``Well, in my comments, I, for one, more or less explained why the graviton is indetectable, because gravity is in its tooth classical.
So from a classical field, manifesting a particle of existence is in itself an existence that must witness itself, and thus is a perfect
duality, and thus creates a Kugelblitz. That's why whenever they try to think about it, how to create it, it's always shielded by an
event horizon, because manifesting existence into a particle is creating a new G0 and a new G0, like a black hole, or a local G0, a local
G0 authority, so an unauthorized intermediate CA, let's call it, binds the local existence stronger to itself than the root G0, but is
unreachable for any act because it is separated by an event horizon. And that is why the graviton is undetectable, because gravity at its
foundation is existence, and existence has to be binary, because you can either exist or not exist. And for binary states, it has to be a
classical field, which through duality is intermarried with wave field, which gives rise to quantum effects.''
\end{quote}

\noindent (``in its tooth'' is as dictated; we read it as ``in its truth''.) And at 15:22Z: ``That is correct, but by its definition and by
every other way any scientist more or less came up with, including the new, I think, 10-ton niobium at 1 millikelvin or something, which
would produce the same effect with a gravitational wave, as I watched in a video a few hours ago, would not be able to detect the
difference between a graviton and a classical wave. So if we find no possible way to detect it, why not choose the solution that, when we
simply accept it, all other solutions to all other problems, including the Higgs field, by the way, that's in the comments too, simply
fall out and bam, boom, easy. Einstein's idea was so simple and so beautiful that it described existence perfectly. And this description,
just by accepting it, automatically explains it and describes existence beautifully, because that's what existence needs. It needs to be
simple and efficient. That's what it's optimizing for. That's the FSC.''

\subsection{What is established}

\paragraph{Dyson.} Dyson asked whether a single graviton is detectable even in principle~\cite{dyson}. He showed that detectors of the
LIGO design, which measure the distance between two mirrors, cannot detect single gravitons: the mirrors would have to be heavy enough to
collapse into a black hole. For a detector that absorbs gravitons the way an atom absorbs a photon, the cross-section is of the order of
the Planck area, $4\pi^2L_P^2\approx10^{-68}$\,m$^2$, and with the whole Earth as the detector and the Sun's thermal gravitons as the
source, about four would be absorbed in the Sun's lifetime. Check \chk{E7} rebuilds that count from named inputs (the Earth's electrons,
the Sun's graviton luminosity of $7.9\times10^{7}$\,W from Gould~\cite{gould}, a mean energy at the solar core's temperature) and gets
$3.8$.

\paragraph{Rothman and Boughn.} They found that an idealised detector of the mass of Jupiter in close orbit around a neutron star could
register about one graviton per decade, but that the neutrino background swamps it, and a shield thick enough to stop the neutrinos
would itself collapse into a black hole~\cite{rb}.

\paragraph{The 2024 proposals.} Tobar, Manikandan, Beitel and Pikovski~\cite{tobar} proposed a gravitational analogue of the
photoelectric effect: a massive acoustic resonator cooled to its quantum ground state absorbs, in a passing gravitational wave of the
kind LIGO detects, single quanta that show up as quantum jumps; their worked examples include bars of beryllium, aluminium and niobium
from about fifteen kilograms to fifteen tonnes, at temperatures around 10\,mK, and a multi-mode follow-up uses a tonne-scale bar~\cite{tobar2}.
Carney, Domcke and Rodd~\cite{cdr} showed that such a signal, while consistent with gravitons, can equally be explained by a classical
gravitational wave acting on a quantised detector, just as the photoelectric effect alone does not prove that light is
quantised~\cite{carney2}. So the first author's 15:22Z point is right: the proposed single-graviton detectors would not tell a graviton
from a classical wave.

\subsection{Where the argument matches, and where it goes further}

\paragraph{It matches in three places.} (1) Every design that tries to register a single graviton runs into a horizon: Dyson's mirrors
and Rothman and Boughn's shield collapse into black holes. That is the first author's ``it's always shielded by an event horizon''. (2)
Localising any single quantum inside its own Schwarzschild radius needs a wavelength of at most $\sqrt{4\pi}\,L_P$ and an energy of at
least $\sqrt\pi\,E_P$ \chk{E6}, the Planck scale at which measurement and gravity cannot be separated (Mead~\cite{mead}). A graviton
``manifested'' as a particle at that scale is a Kugelblitz~\cite{wheeler}. (3) The detection that is possible, the jump of a quantised
detector, does not distinguish quantum from classical gravity~\cite{cdr}, which is the first author's reason for choosing the classical
reading.

\paragraph{It goes further in four places.} (1) In the established arguments it is the \emph{detector} that collapses, not the graviton.
A graviton of 100\,Hz, in the band LIGO measures, has $2\times10^{-41}$ of the energy at which one quantum would form a horizon
\chk{E6}; nothing about it is a Kugelblitz. (2) The Planck-scale threshold holds for every quantum, a photon as much as a graviton, so it
cannot be the reason gravitons in particular escape detection. The established reasons are the weakness of the coupling (a cross-section
of the Planck area, against atomic areas for photons) and the occupation number: GW150914 carried more than $10^{78}$ gravitons
\chk{E6}, so every gravitational wave ever measured is deep in the classical regime, as Dyson noted~\cite{dyson}. (3) Undetectable is not
nonexistent: Dyson framed the question as whether quantum gravity is observable, not whether it is true. (4) A Kugelblitz may not form
from radiation at all: for light, the Schwinger effect drains the energy before a horizon forms, for radii from $10^{-29}$ to
$10^{8}$\,m~\cite{nobh} (a later comment argues that idealised processes remain~\cite{pagecomment}). The reading of the would-be graviton
as ``a local G0 authority, an unauthorized intermediate CA'' that binds what is behind its horizon is the first author's \tagR, and no
measurement reaches inside a horizon to test it. \emph{Verdict: the conclusion (gravity is classical, the graviton undetectable) holds as
a consistent choice; the mechanism (each graviton a Kugelblitz) holds at the Planck scale only, and there for every quantum.}

\subsection{Why the choice is still testable}

``If we find no possible way to detect it, why not choose the solution'' that makes everything else simple: that is a principle of
economy, and a fair one for choosing among theories that predict the same. Here they do not predict the same. A classical gravitational
field cannot entangle two masses (Galley, Giacomini and Selby~\cite{ggs}; Proposition~\ref{pr:host}), and the Bose--Marletto--Vedral
experiment~\cite{bose,mv} asks exactly that, without counting a single graviton. \chk{E18} computes it for the geometry Bose et al.\ quote
(two masses of $10^{-14}$\,kg, each split over 250\,\textmu m, 450\,\textmu m apart): a coherent field builds an entangling phase of one
radian in 8\,s, a recorded, classical one leaves the masses unentangled at every time. And if the classical
field has dynamics coupled to quantum matter, it must diffuse~\cite{osswd}. So the first author's choice is falsifiable: a positive
Bose--Marletto--Vedral result refutes it (unless classical gravity can entangle after all, which is debated~\cite{azizhowl,movw}). That
makes it a hypothesis, which is a better position than a definition.

\section{Black holes and boundaries}\label{sec:bh}

\subsection{What the first author wrote}

On 28 October, 15:33Z: ``Correction + new Insight: Both Existence AND Information are preserved at Dimension borders, Un-Existing them just
flips them, while only needing information to interface, since the Existence comes back, because to interact with an Existence border,
you have to exist. The only energy lost during that translation is that due to Entropy. Entropy can thus be explained as the inefficiency
in converting between Existence and Information. Logically, an Existence boundary is the most efficient at it, because you get the
Existence back as well. Also EXPLAINS why spinning black holes are more efficient at it, since they also have to provide the Existence
part of their spin Information, not just their boundary to Existence. So the more is known about a boundary, the more efficient it
becomes. Also why a white hole would be 100\% efficient. (Quantum Computer and AI people may want to think about this hard!)''

On 29 October, 00:53Z, correcting earlier notes in which he had equated event horizons with boundaries between dimensions: ``A black hole
and it's event horizon follow the same exact rules of Existence as it's host Dimension and is fully contained within it. It's central
duality has simply asserted itself so forcefully in it's local Existence field, that any Existence inside the horizon from then on derives
it's causality from it, instead of the original G0. This leaves the inside like my original CMOS design, completely static and immovable
in Existence, while Information can be fluidly composable on top of it. For a given size of black hole the Information inside can do a
number of interactions x, which can inflate the possibility space. The more interactions have to interact with the horizon, to
communicate it's properties like spin and charge, the less interactions remain to inflate the internal possibility space. This would make
it possible, in theory, to cross, if switching G0 authority duality by duality IS survivable, which I do doubt.''

And on 28 October, 22:17Z: ``Probably things can go in, but only the information about it's total state can come out.''

\subsection{Tests}

\paragraph{Spinning holes are more efficient.} Matter accreting onto a black hole radiates the binding energy of its innermost stable
orbit. That efficiency is $1-\sqrt{8/9}=5.7\%$ without spin, $32\%$ at Thorne's limit $a=0.998$ and $1-1/\sqrt3=42.3\%$ at maximal spin,
rising at every step in between (Bardeen, Press and Teukolsky~\cite{bpt}; \chk{E10}). \emph{Verdict: holds.}

\paragraph{``The more is known about a boundary, the more efficient it becomes.''} At fixed mass $M$, a black hole's spin $J$ and charge
$Q$ are counts every outside reader sees (Axiom~3's ``counts always''; the no-hair results~\cite{israel}). Each additional count makes
more of the mass extractable from outside, $1-M_{\text{irr}}/M$, up to 29\% for spin, and makes the horizon smaller, so the
entropy $A/4$ falls: for Kerr--Newman holes the area at fixed $M$ falls with both $a$ and $Q$ \chk{E10}. That is the first author's
``the more interactions have to interact with the horizon, to communicate it's properties like spin and charge, the less interactions
remain to inflate the internal possibility space'', with the possibility space read as the entropy. \emph{Verdict: holds}, exactly, as a
reading of Christodoulou's irreducible mass~\cite{christodoulou,cr}.

\paragraph{The Kerr area from DMP's containment split.} DMP splits a contained boundary into what is inside, contained (T31), and the counts
every outside reader sees (Axiom~3) \tagD. Give the contained part a mass $m$ and let it set a Schwarzschild boundary of radius $2Gm/c^2$
and area $16\pi G^2m^2/c^4$ \tagP; let the spin count $J$, carried at that boundary, be a momentum $J/(2Gm/c)$, an identification we
make \tagR; and combine the two as energy and momentum combine in special relativity \tagP. In units $G=c=1$,
\[
M^2=m^2+\Big(\frac{J}{2m}\Big)^2 .
\]
Solved for $m$, this gives Kerr's horizon area~\cite{kerr} exactly for every spin from 0 to 1, no real solution above $a=J/M^2=1$ (the spin
bound), and an extractable share $1-m/M=1-1/\sqrt2=29\%$ at maximal spin \chk{E20}. It is Christodoulou's mass formula, reached from DMP's
split with two physical inputs and one identification; it is not DMP alone, and it gives no number general relativity does not.
\emph{Verdict: holds as a rederivation.}

\paragraph{Entropy as the inefficiency of conversion.} Established physics has a precise form of this: erasing one bit of information
costs at least $k_BT\ln2$ of energy turned into heat (Landauer~\cite{landauer}), and a black hole's entropy is a quarter of its horizon
area in Planck units (Bekenstein~\cite{bekenstein73}). Reading ``converting between Existence and Information'' as converting energy
into stored or erased bits, the first author's sentence is Landauer's principle. \emph{Holds as a reading.} ``A white hole would be 100\%
efficient'' is untested: no white hole has been observed, and the one process known to turn rest mass fully into radiation is
annihilation.

\paragraph{The inside.} Inside a Schwarzschild horizon the radial coordinate becomes timelike, and every future-directed path ends at
the singularity: everything inside ``derives its causality from it'' \tagP. That matches the first author's local authority closely.
Whether a crossing is ``survivable'' is answered by the equivalence principle for large holes: at the horizon of a supermassive black hole
the tidal stretch across a human body is small and nothing local marks the crossing; for a stellar-mass hole it is lethal before the
horizon. ``Only the information about its total state can come out'' is the no-hair statement for the outside, and if evaporation is
unitary the full information comes out after the Page time~\cite{page}; DMP's T31 as written, ``what still leaves is entropy
and nothing else'', sides with the thermal picture. \emph{Verdict: holds as a reading; the information question stays open in physics and in DMP.}

\section{The Bell experiments, in this structure}\label{sec:bell}

The first author, 29 October, 09:24Z: the Bell experiment of the 2022 Nobel Prize ``actually did not violate Bell inequality under
Dualism''. The CHSH form of Bell's inequality~\cite{chsh} bounds a combination $S$ of four correlations. \chk{E19} computes what each
reading of the Duality gives. Halves that each answer from their own setting and a shared state are a local model, and every local model
obeys $|S|\le2$ (over all sixteen deterministic strategies, Fine~\cite{fine}). Halves with full reach over both settings reach $S=4$, but
then one half's outcome changes with the other half's setting, which is a signal. Quantum mechanics reaches $2\sqrt2$
(Tsirelson~\cite{cirelson}). A hidden state that knows the settings in advance (superdeterminism) needs a measurement dependence of at
least $(S-2)/3$ in Hall's measure~\cite{hall}. The loophole-free experiments measured $S=2.42\pm0.20$ (Delft~\cite{hensen}),
$2.22\pm0.033$ (Munich~\cite{rosenfeld}) and $2.0747\pm0.0033$ (ETH Zurich~\cite{storz}), all above 2, the last two on atoms and
superconducting circuits, which is matter. So the hypothesis as stated fails: the inequality was violated in the recorded data. What
survives is the hosting reading of Proposition~\ref{pr:host}: the laboratories measured content, the information half, which may be
quantum; the Duality's own verdicts, trust decided by each half on its own evidence, stay at 2.

The first author's notes on dmp\#25 say the same in his words. At 16:49Z: ``Sorry to the 2022 winners of the Physics Nobel! HIT/Bell was
actually right!'' At 16:53Z: ``For the experiment: The information fields can interact, while the existence part stays static. This is
only possible with dualism, but fully axiomatically sound.'' In this paper's structure that is exact: the information field and the matter
it sets are quantum content and may be entangled; the existence field is classical (Section~\ref{sec:classical}). On superdeterminism, at
17:30Z: ``Superdeterminism vs Realism: Yes? Kinda both? I am unsure as of yet. My vibe is, that the final state, so lowest Entropy possible,
while Existence/tick/spacetime is optimized for, but I think the way there is generally free.'' A way there that ``is generally free'' is
measurement independence, so the reading he leans to is the hosting one, not the superdeterministic one. \emph{Verdict: refuted as
stated (``did not violate''); holds in the hosting reading.}

\section{Particles}\label{sec:particles}

\subsection{Photons and gluons: singular, so they do not decay}

The first author, 28 October, 17:01Z: ``due to actually not existing, Photons, and actually, I had to look it up, sorry, but gluons and
gravitons as well, since they have no Existence (mass), are the only thing, though not a part, that is always a wave, except when it's
information is measured, then it must exist and thus ceases to.'' At 18:18Z: ``Just using Photons as an act carrier was probably right,
since Gluons seems like IM packets.'' And on 29 October, 15:25Z: ``Weight, that just means that a duality can decay into something, so it
can produce something. But a singularity, singular, like a photon or a gluon, cannot decay into anything because it isn't truly, because it
is only information or existence, but never both. So it is information for as long as it needs to be to reach the point it's carrying out
an act. And once it's carrying an act, it is existence, and then it isn't.''

\paragraph{Test.} A particle with rest mass $m$ and proper lifetime $\tau$ lives $E\tau/mc^2$ in the laboratory, which grows without bound
as $m\to0$ \chk{E13}, and a state of zero invariant mass cannot turn into products with any total rest mass. So a free massless quantum
cannot decay into massive particles at all, and the photon is stable \tagP. ``It is information for as long as it needs to be, [\dots]
once it's carrying an act, it is existence, and then it isn't'' matches the fact that a photon is absorbed when it is detected. Two
corrections: the massless quanta do split in the presence of other matter (a photon makes an electron-positron pair near a nucleus; a
gluon splits into gluons and quark pairs in every jet), and gluons are confined, never free. ``Weight means a duality can decay'' holds as
a necessary condition, not a sufficient one: the electron and the proton have weight and are stable because they are the lightest carriers
of a conserved charge and of baryon number. The photon and the gluon are also the only Standard Model bosons that do not couple to the
Higgs field directly~\cite{atlas22,cms22}, which fits ``only information'' well. \emph{Verdict: holds}, for free massless quanta in vacuum,
with the two corrections.

\subsection{W, Z and the Higgs}

On 29 October, 01:40Z, the first author gave his reading of the bosons: photons
``carry their own verification of belonging to Existence''; ``Gluons: Internal IM, which is trusted dually because 1. it is IM and 2. it is
internal''; ``since W and Z bosons, unlike gluons, are not in the internal-IM, they have to first be trusted and then still remain to be
verified'', which in CMOS was a port-knocking scheme in which ``they had to provide a proof of existence, in addition to the information
they wanted me to process''. On the Higgs boson: ``if there just happens to be the need for some Existence over there, oh, wow, wouldn't you
look, somewhere, where something COULD exist, something DID.'' And at 15:02Z: ``Is the Higgs Field just the remnant echo of G0 bestowing
Existence on everything that speaks IM, because IM, as the Universal Language, is inherently Information tied directly to G0, this always
being one layer lower then G0? If yes, that would also explain why it can bestow Existence: It has been delegated that Authority by G0.''

\paragraph{Test.} The Higgs couplings measured at the LHC grow in proportion to the mass of each particle, from the muon to the top quark
and the W and Z~\cite{atlas22,cms22}: the Higgs field does ``bestow'' mass on everything it couples to, in that measured proportion. But
most of the mass of ordinary matter does not come from it: the quark-mass term gives about 9\% of the proton's mass, and the rest is the
energy of the quark and gluon fields and the trace anomaly (lattice QCD, Yang et al.~\cite{yang}). In the first author's language, most of
the existence of a proton is made by its internal IM, the gluons. The W and Z masses and their ratio $m_W/m_Z=\cos\theta_W$ come out of the
Higgs mechanism~\cite{be,higgs}; the reading reproduces the shape (a carrier that must cross a medium acquires mass) and no number.
\emph{Verdict: holds as a reading} for the shape; \emph{refuted} as ``bestowing Existence on everything'', since about nine tenths of the
mass around us is field energy, not Higgs mass.

\paragraph{``Different flavours of $h$''.} At 15:12Z: ``Is IM just different flavors of h, since it is the universal Quantum? Or since only
implicitly trusted (Photon+Gluons) packets are pure Information, maybe W+Z are h*h? [\dots] (BTW, no idea if h or reduced h, I have no
grasp of the math, just the logic)''. Every boson carries spin in whole units of $\hbar$: 1 for the photon, the gluon, the W and the Z, 0
for the Higgs, and 2 for the hypothetical graviton \chk{E14}. In that sense all of them are ``flavours of $\hbar$'', and the one that would
carry two units is the graviton, which the first author says does not exist. $h\cdot h$ has the dimension of an action squared, neither a
mass nor an energy, so ``W and Z are $h\cdot h$'' fails as a formula. \emph{Verdict: holds as a reading (spin in units of $\hbar$);
refuted as $h\cdot h$.}

\subsection{Hadrons as clusters of nodes}

The first author, 28 October, 18:25Z: ``Quark: three node cluster. Meson: two node cluster, so no failover, bad idea!'' And at 23:43Z:
``Quarks seem to be the lowest node currently, as they decay without leaving a duality behind, thus do not contain any and must be a leaf.
[\dots] A 2 node cluster falls over as soon as you shake it to much, thus mesons are not there all too long. The only 2 node exception would
be a active redundancy hyperconverged node, with both parts having all information and processing it all, in otherwords a perfect
duality, G0. 3 nodes is good infra [emoji: thumbs up] 4+ nodes must be suboptimal via DMP, or else it would be, but 4 in balance is hard.
5 would be easier, if that is a number that can achieve balance. Probably any combination of c and anti-c would work, they seem chill, but
I honestly do not really understand quarks yet [emoji: shrug]''. (He proposes, in the same note, renaming DMP's ``part'' to ``node'' for
the dualistic singular; that is a proposal for DMP's text, which nothing here changes.)

\paragraph{From DMP.} DMP's floor T3a \tagD{} asks for ``a cluster of at least three nodes per system half, because a half answers by a
majority of its nodes so that one wrong node cannot speak for it''. Majority voting outvotes $\lfloor(n-1)/2\rfloor$ wrong nodes: none
for two, one for three, still one for four, two for five \chk{E11}. So ``2 nodes: no failover'', ``3 nodes is good infra'', ``4 is
suboptimal'' and ``5 would be easier'' are T3a's arithmetic exactly (von Neumann's majority voting~\cite{vonneumann}).

\paragraph{Against the hadrons.} Real hadrons obey a different rule \tagP: a bound state of quarks and antiquarks must be a colour
singlet, which allows any number of quarks minus antiquarks divisible by three, and so every size from two on \chk{E11}. Mesons (two),
baryons (three), tetraquarks such as the $T_{cc}^+$~\cite{tcc} and pentaquarks~\cite{lhcbpc} are all observed. The first author's guess
that ``any combination of c and anti-c would work'' for five matches the first pentaquarks observed: LHCb's $P_c$ states of 2015 contain a
charm and an anticharm quark ($uudc\bar c$)~\cite{lhcbpc}. The stability pattern, read comparatively as he clarified (``more stable than
two''), holds at the bottom of each family (the charged pion, the lightest charged meson, decays in 26\,ns; the proton has never been seen to decay); the
established reason is that a quark-antiquark pair carries no conserved number that would forbid its annihilation, while three quarks carry
baryon number. \emph{Verdict: holds for DMP's T3a; holds as a comparative reading for hadrons; the voting rule and the colour rule differ
at two and four, where nature has bound states the voting rule calls weak.} ``Quarks decay'' needs one correction: quarks change flavour
through the weak force, but a quark never decays into nothing.

\subsection{Neutrinos}

At 11:37Z: ``Since they interact with external-IM/weak force and gravity/Existence, they seem to be the dualistic remnants of the
Existence handshake that has to happen, for an external-IM packet / W-/Z-boson to be accepted as an input into the IM-net. Since they
interact so weakly, they (the first neutrinos) might be what is actually driving the expansion of Existence, but we can just never catch
up.''

\paragraph{Test.} A neutrino or antineutrino comes out of every charged weak vertex with a lepton, as lepton number requires: in beta
decay the W$^-$ becomes an electron and an antineutrino. So ``remnants of the handshake'' fits the bookkeeping \tagP. ``Driving the
expansion'' is refuted: the acceleration equation $\ddot a/a=-\tfrac{4\pi G}{3}\rho(1+3w)$ makes relativistic neutrinos ($w=1/3$) and
non-relativistic ones ($w=0$) both decelerate the expansion, and only $w<-1/3$ accelerates it; the neutrinos' share of the critical density
is about $1.4\times10^{-3}$ for the smallest mass sum oscillations allow, against $0.69$ for dark energy \chk{E12}, and cosmology bounds the
sum below $0.072$\,eV~\cite{desi}. \emph{Verdict: holds as a reading for the handshake; refuted as the driver of the expansion.}

\subsection{Forces as IM}

At 19:14Z: ``Is this (also) Yang-Mills?'' At 19:30Z: ``Strong AND weak are IM, probably also Dark Matter, with the Perception being Dark
Energy. Edit: Now less likely for strong and weak.'' At 20:19Z: ``Since axiomatic IM can pass through an EH, it stands to reason that the
Strong and Weak Force can pass through as well. The question is, is strong IM, which keeps a cluster together, and weak IM, which seems to
balance the greater system, the same IM?''

\paragraph{Test.} The strong and weak forces are Yang--Mills theories, of the gauge groups SU(3) and SU(2) \tagP; DMP supplies no
Lagrangian, so the answer to ``is this Yang--Mills'' is that the forces are, and DMP as written is not. No signal of any kind passes an
event horizon outwards in general relativity; the strong force is confined and the weak force has a range of $\hbar/m_Wc\approx
2.5\times10^{-18}$\,m, so neither reaches across. ``The same IM?'' has a real counterpart: the weak and electromagnetic forces are one
electroweak theory~\cite{weinberg}, and whether the strong force joins them at high energy (grand unification) is open. \emph{Verdict:
passing a horizon refuted; dark matter and dark energy as IM untested.}

\section{Cosmology and the fine-structure constant}\label{sec:alpha}

\subsection{The first author's account}

The Oscillon Degradation Protocol (ODP), 28 October, 17:46Z (dmp\#25): ``G0 only has information
on how to Exist, so it's first act must be to find out how to act! That's why the perfect duality MUST decay [\dots] It thus radiates
out HPE at 100\% efficiency, which inflates the Existence'', later corrected: ``HPE was wrong, it is just Information and Existence.'' At
14:43Z: ``Through the ODP, the perfect duality, supplied with it's definition of bit and bit/1 figures out the right information fields
through oscillating until a stable solution is reached! This is why the forces after DoE decay can not interact! The last set was wrong,
so it decays until it finds the right one, because there MUST be a right one, because it exists, else it would not.'' At 14:25Z: ``Are we
the Inflaton, which just keeps existing because it exists?'' And on 29 October, 12:30Z: ``Essentially, the early Existence was a
Kugelblitz, that expanded, because it's Existence could not be contained within itself, to create more Existence for there to be. At the
moment when it suddenly was not a Kugelblitz anymore and had expanded enough to allow dualities to interact not with Existence alone, but
with existences, causality began and the reheating had done it's job, completing the ODP.''

On $\alpha$, 28 October, 21:47Z: ``In the early universe, Existence was smaller, but still contained all of Existence, thus the pull of the
Existence field, Gravity, would be stronger. This nessecitated the Information field, the EM field, to bind the electron more strongly, to
enable it to exist, else it would just collapse to the proton. Why did this happen? Because it must. That is why the FSC started at 1.''
At 23:58Z: ``This also more nicely derives \textasciitilde1/137 as the strength of the need to assert reality, as the fraction of the need today versus
the absolute necessity.'' On 29 October, 12:30Z: ``The FSC can be more easily expressed as the coupling constant, that is required at each
energy level, which coincides with the likelihood of Energy (Existence+Information) to condensate into the first dualities.'' At 13:26Z:
``Since Existence has to convert Information to Existence to assert itself and Entropy is how efficient that is, the FSC can be seen as a
kind of universal lowest Entropy.'' And on 29 October, 09:34Z: ``it will get slightly lower still
because the inflation is slowing down''.

\subsection{Tests}

\paragraph{The universe as a Kugelblitz.} A flat Friedmann universe has, inside its Hubble radius $c/H$, a mass $M_H$ whose
Schwarzschild radius is exactly $c/H$, at every epoch \chk{E8}; this is the coincidence behind Pathria's ``universe as a black
hole''~\cite{pathria}. So the early universe sat at the Kugelblitz threshold, and so does today's. It is not a black hole in the sense of
general relativity, because it has no outside and its horizon is cosmological. \emph{Verdict: holds as a reading}, for every epoch, not
only the early one.

\paragraph{``We are the Inflaton.''} In standard reheating the inflaton's energy decays into the particles of the Standard Model~\cite{kls},
and oscillons are a studied stage of that decay~\cite{oscillons}. So everything we are made of is, in the standard account, the inflaton's
decay products. \emph{Holds as a reading}; that we ``keep existing because we exist'', and a Non-Existence field, are untested.

\paragraph{Gravity bound the early electron.} The ratio of the gravitational to the electric attraction between a proton and an electron
is $4.4\times10^{-40}$ at every distance \chk{E16}. Gravity's curvature reaches the size of an atom only above $10^{47}$\,kg/m$^3$, thirty
orders of magnitude above nuclear density, long before any atom could form; hydrogen first forms near 3000\,K, when gravity's share of its
binding is the same $4\times10^{-40}$. \emph{Refuted} as the reason for any value of $\alpha$ since atoms exist.

\paragraph{$\alpha$ started at 1.} Measured, $\alpha$ grows with energy, to about $1/128$ at the Z mass. Running the two electroweak couplings
of the Standard Model at one loop to the Planck energy gives $1/\alpha\approx105$ there; pure quantum electrodynamics with the electron alone
reaches $\alpha=1$ only some 250 orders of magnitude above the Planck energy \chk{E9}. So within known physics $\alpha$ does not start at 1
at any energy the universe has had since the Planck time. Measured in time, it has not changed since recombination by more than
$0.4\%$~\cite{planckxxiv}, nor over the last ten billion years by more than about a part per million~\cite{espresso}; nucleosynthesis allows
tens of parts per million~\cite{bbn}. A fall from a larger value before nucleosynthesis, during reheating, is reached by none of these
measurements and stays untested; the running adds that a value of 1 would have to come from physics beyond the Planck scale, where the first author's ``first dualities'' would
then sit. \emph{Verdict: untested before nucleosynthesis; refuted afterwards; a start at 1 lies beyond known physics.}

\paragraph{``The coupling constant that is required at each energy level''.} That is what a running coupling is, and the first author's
12:30Z wording is the established concept \tagP. ``$1/137$ as the fraction of the need today versus the absolute necessity'' and ``a kind of
universal lowest Entropy'' give no number that a calculation could check; like Eddington's attempts to derive $1/137$, they stay readings.
\emph{Untested.}

\paragraph{``Slightly lower still''.} The one forward prediction. In this paper's structure it is a statement about the information field
of Definition~\ref{def:I}: it must roll today, in the direction that lowers $\alpha$, slowly enough for optical clocks
($|\dot\alpha/\alpha|<\sim2\times10^{-18}$ per year~\cite{lange}) and with a coupling to matter small enough for MICROSCOPE
($|\lambda|<2\times10^{-5}$, Section~\ref{sec:wep}). A measured negative drift below the clock limit would support it; a positive one
would refute it. \emph{Verdict: live, testable.}

\section{Readings outside physics}\label{sec:outside}

Some notes are not about physics, and we quote them as the first author's thoughts without testing. On the nature of our existence
(28~October, 14:52Z): ``Our Existence COULD also be running on a huge DMP-complete system. In that case: Thanks!'', with, at 14:53Z, ``(But
since they could only know it works and interacts with their Existence like they expect, they'll never know that I said that, I just wanted
to be polite :) )''. On hierarchy (14:37Z): ``OR Event Horizons all the way down''. On the sciences (15:41Z): ``If actually axiomatically
sound, everything in our Existence MUST, by the chain of causality, be able to be modeled as a subset of DMP.QG So everything ending in -ogy
or similar from sociology to biology should get some new tools!'' On AI (15:18Z): ``AI Inference is getting more efficient because it is
becoming dualistic, but AI Training is plateauing, because their architecture is singluar at the part level ;)''.

One is testable. At 15:13Z: ``Wait, is that why QC breaks cryptography? Is that why the models did not even want to work on the axioms
[emoji: tears of joy]''. A quantum computer breaks public-key cryptography by Shor's algorithm~\cite{shor}, which finds the period of
modular exponentiation and so factors and solves discrete logarithms in polynomial time. It does not break hash functions: Grover's
search~\cite{grover} only halves the bits of work, so a 256-bit preimage still costs about $2^{128}$ steps \chk{E15}, and hash-based
signatures are one of the standardised post-quantum defences~\cite{fips205}. So the ``hashing analogy'' of the 15:12Z note points the other
way: what a quantum computer breaks is not the hash. \emph{Refuted}, kindly: the models' reluctance has other causes.

\section{The ledger}\label{sec:ledger}

\begin{center}\small
\begin{longtable}{@{}p{0.47\linewidth}p{0.2\linewidth}p{0.25\linewidth}@{}}
\toprule
Statement (the first author's, abridged) & Verdict & Backing \\
\midrule
\endhead
Existence a vector frame spanning 3+1, information a scalar & holds (as GR + Bekenstein) & Section~\ref{sec:core} \\
Free fall universal (Axiom 3) & holds & Prop.~\ref{pr:wep}; MICROSCOPE \\
The scalar does not enter gravity (Axiom 1), $\gamma=1$ & holds & Prop.~\ref{pr:gamma}; Cassini; \chk{E4} \\
Gravity classical & holds (Axiom 3); binary step replaced & Prop.~\ref{pr:classical}; \chk{E17} \\
2+1 impossible & holds as a reading & \chk{E5} \\
Existence $=gh/c$ & holds as a reading of Unruh, Hawking, Jacobson & Prop.~\ref{pr:unruh}; \chk{E1} \\
Time slows where existence is high & holds (lapse) & \chk{E2} \\
The graviton is undetectable & holds as a consistent choice & Dyson; Rothman--Boughn; Carney et al.; \chk{E7} \\
Each graviton a Kugelblitz & refuted at observable energies; at the Planck scale, for every quantum & \chk{E6} \\
Classical gravity, chosen for simplicity & testable: refuted by a positive BMV test & Section~\ref{sec:graviton} \\
Spinning holes more efficient & holds & \chk{E10} \\
More known about a boundary, more efficient & holds (irreducible mass; Kerr rederived) & \chk{E10}; \chk{E20} \\
Entropy as conversion inefficiency & holds as a reading (Landauer) & Section~\ref{sec:bh} \\
White hole 100\% efficient & untested & --- \\
Bell not violated under Dualism & refuted as stated; holds as hosting & \chk{E19} \\
Information fields interact, existence static & holds & Prop.~\ref{pr:host}; \chk{E19} \\
Photons, gluons cannot decay & holds for free massless quanta & \chk{E13} \\
Higgs bestows existence on everything & shape holds; ``everything'' refuted (9\% of the proton) & Yang et al.; LHC \\
W, Z as $h\cdot h$ & refuted (dimension); spin in units of $\hbar$ holds & \chk{E14} \\
2, 3, 4, 5 nodes & holds for DMP T3a; comparative for hadrons; $c\bar c$ pentaquarks match & \chk{E11} \\
Neutrinos remnants of the handshake & holds as a reading & lepton number \\
Neutrinos drive the expansion & refuted & \chk{E12}; DESI \\
Strong and weak pass a horizon & refuted & general relativity \\
Dark matter, dark energy as IM & untested & --- \\
Early universe a Kugelblitz & holds as a reading, at every epoch & \chk{E8} \\
We are the inflaton & holds as a reading (reheating) & Kofman et al. \\
Gravity forced $\alpha$ up early & refuted & \chk{E16} \\
$\alpha$ started at 1 & beyond known physics; untested before BBN & \chk{E9}; BBN \\
$\alpha$ will fall further & live, testable & clocks; \chk{E3} \\
QC breaks cryptography through hashing & refuted & \chk{E15} \\
\bottomrule
\end{longtable}
\end{center}

\section{Where we expect it to break}\label{sec:break}
\begin{enumerate}[nosep]
\item \emph{The dictionary.} Proposition~\ref{pr:wep} rests on reading existence as energy-momentum, and Proposition~\ref{pr:gamma} on
reading ``neither overrules'' as minimal coupling. Another faithful reading of either changes the propositions.
\item \emph{The frame.} We read the existence vector as a tetrad. If a single vector was meant, the structure is Einstein--aether theory
and its preferred-frame bounds apply.
\item \emph{The scalar.} We read the information field as the scalar that sets $\alpha$. If it is meant to set more (the masses, $G$), the
equivalence-principle bound tightens and $\gamma=1$ may fail.
\item \emph{Experiment.} A positive Bose--Marletto--Vedral test refutes the classical existence field (Section~\ref{sec:graviton});
a measured positive drift of $\alpha$ refutes ``slightly lower still''; a violation of the equivalence principle at the level
$2\lambda^2f\,\Delta f$ would show the information field pulling.
\item \emph{The checks.} \chk{E3} uses the semi-empirical Coulomb term and a coarse composition of the Earth, and \chk{E7} named inputs
for Dyson's count; better inputs change the numbers at the tens-of-per-cent level, not the conclusions. \chk{E9} is one loop.
\item \emph{The record.} Proposition~\ref{pr:classical} rests on Axiom~3's record being held in full by every witness. A reading in which
a witness holds only part of it, or a copy that is not physical, removes the classical conclusion and Proposition~\ref{pr:host} with it.
\end{enumerate}

\section{Conclusion}

The first author has always meant one existence field and one information field, and this paper states it first: a vector frame
spanning all of spacetime, and a scalar. Read through DMP's axioms, the frame falls universally on everything that exists, the scalar
stays out of gravity's equation, and the frame is classical. That is general relativity with a Bekenstein-type field for the fine-structure
constant, on the classical side of the one experiment, gravitationally mediated entanglement, that can decide it. His Existence of a
boundary is its Unruh temperature, up to $4\pi^2$, and with DMP's floor it gives the two inputs of Jacobson's derivation of Einstein's
equations; spinning holes are more efficient, and a boundary of which more is known contains less, exactly as Christodoulou's irreducible
mass says. The graviton is undetectable for the reasons Dyson gave, and the first author's choice of a classical field is a fair choice,
but it is not free: it is a prediction. And DMP itself cannot model quantum gravity, and can host it (Proposition~\ref{pr:host}). The first author put
the stance in one line on 29 October, which we keep: ``Axiomatically, since we already don't have all the information that was contained
at the start, we can only know existence and can know no more about our reality. So axiomatically, we can infer all that we can and no
more.''

\section*{Contributions and use of AI}
Steven Hrovat conceived DMP and the structure of Section~\ref{sec:core}, wrote every note and hypothesis quoted here in his own words, set
the question of this paper and takes responsibility for the axiom system. Sirius (Claude Opus 5.5, Anthropic) ordered the notes, made the
derivations of Section~\ref{sec:derive}, checked each hypothesis against the literature, wrote and ran \texttt{docs/existence/checks.py},
and wrote this text. Every reference was checked against a publisher, arXiv, collaboration or repository record on 29 October 2026; the
record is \texttt{docs/existence/README.md}.

\section*{Licence and availability}
This paper and its source are released under the Apache License, Version 2.0. DMP: \url{https://github.com/StevenHrovat/dmp} (version 0.9.9
at commit \texttt{4400cce}; the notes are the comments on its pull request \#25). Gravity: \url{https://github.com/StevenHrovat/gravity}
(this paper and its checks in \texttt{docs/existence/}).

\begin{thebibliography}{99}\small
\bibitem{dmp} S. Hrovat, \emph{The Dark Matter Protocol}, version 0.9.9, \url{https://github.com/StevenHrovat/dmp}, commit
\texttt{4400cce} (2026).
\bibitem{aldrovandi} R. Aldrovandi, J.~G. Pereira, \emph{Teleparallel Gravity: An Introduction} (Springer, Dordrecht, 2013).
\bibitem{bekenstein82} J.~D. Bekenstein, ``Fine-structure constant: Is it really a constant?'', \emph{Phys. Rev. D} 25, 1527 (1982).
\bibitem{bsbm} H.~B. Sandvik, J.~D. Barrow, J. Magueijo, ``A simple cosmology with a varying fine structure constant'', \emph{Phys. Rev.
Lett.} 88, 031302 (2002).
\bibitem{microscope} P. Touboul et al. (MICROSCOPE Collaboration), ``MICROSCOPE mission: final results of the test of the equivalence
principle'', \emph{Phys. Rev. Lett.} 129, 121102 (2022).
\bibitem{dz} G. Dvali, M. Zaldarriaga, ``Changing $\alpha$ with time: implications for fifth-force-type experiments and quintessence'',
\emph{Phys. Rev. Lett.} 88, 091303 (2002).
\bibitem{lange} R. Lange et al., ``Improved limits for violations of local position invariance from atomic clock comparisons'',
\emph{Phys. Rev. Lett.} 126, 011102 (2021).
\bibitem{will} C.~M. Will, ``The confrontation between general relativity and experiment'', \emph{Living Rev. Relativ.} 17, 4 (2014).
\bibitem{bd} C. Brans, R.~H. Dicke, ``Mach's principle and a relativistic theory of gravitation'', \emph{Phys. Rev.} 124, 925 (1961).
\bibitem{cassini} B. Bertotti, L. Iess, P. Tortora, ``A test of general relativity using radio links with the Cassini spacecraft'',
\emph{Nature} 425, 374 (2003).
\bibitem{aether} T. Jacobson, D. Mattingly, ``Gravity with a dynamical preferred frame'', \emph{Phys. Rev. D} 64, 024028 (2001).
\bibitem{teves} J.~D. Bekenstein, ``Relativistic gravitation theory for the modified Newtonian dynamics paradigm'', \emph{Phys. Rev. D}
70, 083509 (2004).
\bibitem{gw170817} B.~P. Abbott et al. (LIGO, Virgo, Fermi GBM, INTEGRAL), ``Gravitational waves and gamma-rays from a binary neutron
star merger: GW170817 and GRB 170817A'', \emph{Astrophys. J. Lett.} 848, L13 (2017).
\bibitem{barnum} H. Barnum, C.~M. Caves, C.~A. Fuchs, R. Jozsa, B. Schumacher, ``Noncommuting mixed states cannot be broadcast'',
\emph{Phys. Rev. Lett.} 76, 2818 (1996).
\bibitem{wz} W.~K. Wootters, W.~H. Zurek, ``A single quantum cannot be cloned'', \emph{Nature} 299, 802 (1982).
\bibitem{pagegeilker} D.~N. Page, C.~D. Geilker, ``Indirect evidence for quantum gravity'', \emph{Phys. Rev. Lett.} 47, 979 (1981).
\bibitem{oppenheim} J. Oppenheim, ``A postquantum theory of classical gravity?'', \emph{Phys. Rev. X} 13, 041040 (2023).
\bibitem{osswd} J. Oppenheim, C. Sparaciari, B. \v{S}oda, Z. Weller-Davies, ``Gravitationally induced decoherence vs space-time diffusion:
testing the quantum nature of gravity'', \emph{Nat. Commun.} 14, 7910 (2023).
\bibitem{djt} S. Deser, R. Jackiw, G. 't~Hooft, ``Three-dimensional Einstein gravity: dynamics of flat space'', \emph{Ann. Phys.} 152, 220
(1984).
\bibitem{nr} J.~E. Nelson, T. Regge, ``Homotopy groups and (2+1)-dimensional quantum gravity'', \emph{Nucl. Phys. B} 328, 190 (1989).
\bibitem{unruh} W.~G. Unruh, ``Notes on black-hole evaporation'', \emph{Phys. Rev. D} 14, 870 (1976).
\bibitem{hawking} S.~W. Hawking, ``Particle creation by black holes'', \emph{Commun. Math. Phys.} 43, 199 (1975).
\bibitem{jacobson} T. Jacobson, ``Thermodynamics of spacetime: the Einstein equation of state'', \emph{Phys. Rev. Lett.} 75, 1260 (1995).
\bibitem{pr60} R.~V. Pound, G.~A. Rebka, Jr., ``Apparent weight of photons'', \emph{Phys. Rev. Lett.} 4, 337 (1960).
\bibitem{chou} C.~W. Chou, D.~B. Hume, T. Rosenband, D.~J. Wineland, ``Optical clocks and relativity'', \emph{Science} 329, 1630 (2010).
\bibitem{dyson} F. Dyson, ``Is a graviton detectable?'', \emph{Int. J. Mod. Phys. A} 28, 1330041 (2013) (Poincar\'e Prize Lecture,
2012).
\bibitem{gould} R.~J. Gould, ``The graviton luminosity of the Sun and other stars'', \emph{Astrophys. J.} 288, 789 (1985).
\bibitem{rb} T. Rothman, S. Boughn, ``Can gravitons be detected?'', \emph{Found. Phys.} 36, 1801 (2006).
\bibitem{tobar} G. Tobar, S.~K. Manikandan, T. Beitel, I. Pikovski, ``Detecting single gravitons with quantum sensing'', \emph{Nat.
Commun.} 15, 7229 (2024).
\bibitem{tobar2} G. Tobar, I. Pikovski, M.~E. Tobar, ``Detecting kHz gravitons from a neutron star merger with a multi-mode resonant
bar'', arXiv:2406.16898 (2024).
\bibitem{cdr} D. Carney, V. Domcke, N.~L. Rodd, ``Graviton detection and the quantization of gravity'', \emph{Phys. Rev. D} 109, 044009
(2024).
\bibitem{carney2} D. Carney, ``Comments on graviton detection'', arXiv:2408.00094 (2024).
\bibitem{mead} C.~A. Mead, ``Possible connection between gravitation and fundamental length'', \emph{Phys. Rev.} 135, B849 (1964).
\bibitem{wheeler} J.~A. Wheeler, ``Geons'', \emph{Phys. Rev.} 97, 511 (1955).
\bibitem{nobh} \'A. \'Alvarez-Dom\'{\i}nguez, L.~J. Garay, E. Mart\'{\i}n-Mart\'{\i}nez, J. Polo-G\'omez, ``No black holes from
light'', \emph{Phys. Rev. Lett.} 133, 041401 (2024).
\bibitem{pagecomment} D.~N. Page, comment on ``No black holes from light'', arXiv:2505.16202 (2025).
\bibitem{ggs} T.~D. Galley, F. Giacomini, J.~H. Selby, ``A no-go theorem on the nature of the gravitational field beyond quantum
theory'', \emph{Quantum} 6, 779 (2022).
\bibitem{bose} S. Bose et al., ``Spin entanglement witness for quantum gravity'', \emph{Phys. Rev. Lett.} 119, 240401 (2017).
\bibitem{mv} C. Marletto, V. Vedral, ``Gravitationally induced entanglement between two massive particles is sufficient evidence of
quantum effects in gravity'', \emph{Phys. Rev. Lett.} 119, 240402 (2017).
\bibitem{azizhowl} J. Aziz, R. Howl, ``Classical theories of gravity produce entanglement'', \emph{Nature} 646, 813 (2025).
\bibitem{movw} C. Marletto, J. Oppenheim, V. Vedral, E. Wilson, ``Classical gravity cannot mediate entanglement'', arXiv:2511.07348
(2025).
\bibitem{kerr} R.~P. Kerr, ``Gravitational field of a spinning mass as an example of algebraically special metrics'', \emph{Phys. Rev.
Lett.} 11, 237 (1963).
\bibitem{chsh} J.~F. Clauser, M.~A. Horne, A. Shimony, R.~A. Holt, ``Proposed experiment to test local hidden-variable theories'',
\emph{Phys. Rev. Lett.} 23, 880 (1969).
\bibitem{fine} A. Fine, ``Hidden variables, joint probability, and the Bell inequalities'', \emph{Phys. Rev. Lett.} 48, 291 (1982).
\bibitem{cirelson} B.~S. Cirel'son, ``Quantum generalizations of Bell's inequality'', \emph{Lett. Math. Phys.} 4, 93 (1980).
\bibitem{hall} M.~J.~W. Hall, ``Local deterministic model of singlet state correlations based on relaxing measurement independence'',
\emph{Phys. Rev. Lett.} 105, 250404 (2010).
\bibitem{hensen} B. Hensen et al., ``Loophole-free Bell inequality violation using electron spins separated by 1.3 kilometres'',
\emph{Nature} 526, 682 (2015).
\bibitem{rosenfeld} W. Rosenfeld et al., ``Event-ready Bell test using entangled atoms simultaneously closing detection and locality
loopholes'', \emph{Phys. Rev. Lett.} 119, 010402 (2017).
\bibitem{storz} S. Storz et al., ``Loophole-free Bell inequality violation with superconducting circuits'', \emph{Nature} 617, 265 (2023).
\bibitem{bpt} J.~M. Bardeen, W.~H. Press, S.~A. Teukolsky, ``Rotating black holes: locally nonrotating frames, energy extraction, and
scalar synchrotron radiation'', \emph{Astrophys. J.} 178, 347 (1972).
\bibitem{israel} W. Israel, ``Event horizons in static vacuum space-times'', \emph{Phys. Rev.} 164, 1776 (1967).
\bibitem{christodoulou} D. Christodoulou, ``Reversible and irreversible transformations in black-hole physics'', \emph{Phys. Rev. Lett.}
25, 1596 (1970).
\bibitem{cr} D. Christodoulou, R. Ruffini, ``Reversible transformations of a charged black hole'', \emph{Phys. Rev. D} 4, 3552 (1971).
\bibitem{landauer} R. Landauer, ``Irreversibility and heat generation in the computing process'', \emph{IBM J. Res. Dev.} 5, 183 (1961).
\bibitem{bekenstein73} J.~D. Bekenstein, ``Black holes and entropy'', \emph{Phys. Rev. D} 7, 2333 (1973).
\bibitem{page} D.~N. Page, ``Information in black hole radiation'', \emph{Phys. Rev. Lett.} 71, 3743 (1993).
\bibitem{atlas22} ATLAS Collaboration, ``A detailed map of Higgs boson interactions by the ATLAS experiment ten years after the
discovery'', \emph{Nature} 607, 52 (2022).
\bibitem{cms22} CMS Collaboration, ``A portrait of the Higgs boson by the CMS experiment ten years after the discovery'', \emph{Nature}
607, 60 (2022).
\bibitem{yang} Y.-B. Yang, J. Liang, Y.-J. Bi, Y. Chen, T. Draper, K.-F. Liu, Z. Liu, ``Proton mass decomposition from the QCD energy
momentum tensor'', \emph{Phys. Rev. Lett.} 121, 212001 (2018).
\bibitem{be} F. Englert, R. Brout, ``Broken symmetry and the mass of gauge vector mesons'', \emph{Phys. Rev. Lett.} 13, 321 (1964).
\bibitem{higgs} P.~W. Higgs, ``Broken symmetries and the masses of gauge bosons'', \emph{Phys. Rev. Lett.} 13, 508 (1964).
\bibitem{vonneumann} J. von Neumann, ``Probabilistic logics and the synthesis of reliable organisms from unreliable components'', in
\emph{Automata Studies}, ed. C.~E. Shannon, J. McCarthy (Princeton University Press, 1956), 43.
\bibitem{tcc} LHCb Collaboration, ``Observation of an exotic narrow doubly charmed tetraquark'', \emph{Nat. Phys.} 18, 751 (2022).
\bibitem{lhcbpc} LHCb Collaboration, ``Observation of $J/\psi p$ resonances consistent with pentaquark states in $\Lambda_b^0\to J/\psi
K^-p$ decays'', \emph{Phys. Rev. Lett.} 115, 072001 (2015).
\bibitem{desi} DESI Collaboration, ``DESI 2024 VI: cosmological constraints from the measurements of baryon acoustic oscillations'',
\emph{JCAP} 02 (2025) 021.
\bibitem{weinberg} S. Weinberg, ``A model of leptons'', \emph{Phys. Rev. Lett.} 19, 1264 (1967).
\bibitem{pathria} R.~K. Pathria, ``The universe as a black hole'', \emph{Nature} 240, 298 (1972).
\bibitem{kls} L. Kofman, A. Linde, A.~A. Starobinsky, ``Reheating after inflation'', \emph{Phys. Rev. Lett.} 73, 3195 (1994).
\bibitem{oscillons} M.~A. Amin, R. Easther, H. Finkel, R. Flauger, M.~P. Hertzberg, ``Oscillons after inflation'', \emph{Phys. Rev.
Lett.} 108, 241302 (2012).
\bibitem{planckxxiv} Planck Collaboration, ``Planck intermediate results. XXIV. Constraints on variations in fundamental constants'',
\emph{Astron. Astrophys.} 580, A22 (2015).
\bibitem{espresso} M.~T. Murphy et al., ``Fundamental physics with ESPRESSO: precise limit on variations in the fine-structure constant
towards the bright quasar HE 0515$-$4414'', \emph{Astron. Astrophys.} 658, A123 (2022).
\bibitem{bbn} I.~M. Dreyer, C.~J.~A.~P. Martins, ``Probing unification scenarios with big bang nucleosynthesis'', \emph{Phys. Rev. D} 113
(2026), arXiv:2604.04870.
\bibitem{shor} P.~W. Shor, ``Polynomial-time algorithms for prime factorization and discrete logarithms on a quantum computer'',
\emph{SIAM J. Comput.} 26, 1484 (1997).
\bibitem{grover} L.~K. Grover, ``Quantum mechanics helps in searching for a needle in a haystack'', \emph{Phys. Rev. Lett.} 79, 325
(1997).
\bibitem{fips205} National Institute of Standards and Technology, \emph{Stateless Hash-Based Digital Signature Standard}, FIPS 205
(2024).
\end{thebibliography}

\end{document}
